% Paper 5, §6 — Agentic benchmarks: realism vs measurability
% Anchors: kb/notes/evaluation/agentic-benchmarks.md
%          kb/excerpts/{jimenez2024-swebench,xie2024-osworld,
%                       mialon2023-gaia,yao2024-tau-bench}.md

\section{Agentic benchmarks: what an LLM can DO}
\label{sec:agentic-benchmarks}

The \glsterm{mmlu}{MMLU} lineage (\cref{sec:knowledge-benchmarks}) measures what a model
\emph{knows}: a single text response, scored against a fixed answer \glsterm{key}{key}
by exact match. The lineage taken up in this section measures what a
model \emph{does}: a sequence of actions in an environment --- code
edits, mouse clicks, web fetches, tool calls --- scored against the
\emph{end-state} those actions produce. The thesis is that the four
canonical 2023--2024 agentic benchmarks --- SWE-Bench
\citep{jimenez2024-swebench}, OSWorld \citep{xie2024-osworld}, GAIA
\citep{mialon2023-gaia}, and $\tau$-Bench \citep{yao2024-tau-bench}
--- each pick a different point on a single tradeoff curve: the more
the environment looks like real deployment, the harder it is to score
reliably. SWE-Bench buys ecological realism with the cost of public
GitHub-issue contamination and long-tail spec ambiguity. OSWorld buys
desktop-OS realism with a vision-grounding bottleneck that confounds
the metric. GAIA buys low contamination with hand-crafted items and
a string-match verifier on the final answer field. $\tau$-Bench buys
dialogue-and-policy realism by replacing the user with another \glsterm{llm}{LLM},
which is the easiest piece to simulate and the noisiest one to
interpret. None of the four solves the realism-vs-measurability
problem; each picks a corner of it.

The shared formal frame is the agentic-benchmark instance
$\mathcal{B} = (s_0, \mathcal{A}, \mathcal{T}, \rho, g, V)$%
\footnote{KB note: \texttt{kb/notes/evaluation/agentic-benchmarks.md\#sec-1}.}
--- initial environment state $s_0$, action space $\mathcal{A}$,
transition $\mathcal{T}$, task description $\rho$, goal $g$, and an
\glsterm{end-state-verifier}{end-state verifier} $V(s_T, g) \in \{0, 1\}$ that accepts or rejects
the trajectory. The agent under test maps $(\rho, s_t, h_t) \to a_t$
at each step, with history $h_t$, until a stop condition or step
budget is hit. The structural break from QA benchmarks is that $V$
checks the post-trajectory state, not the trajectory itself: any path
that produces a passing end-state is correct, even ones the
benchmark author did not anticipate.

\subsection{SWE-Bench: software engineering as benchmark}
\label{sec:agentic-swebench}

\citet{jimenez2024-swebench} construct SWE-Bench from real
GitHub-issue $\to$ pull-request pairs%
\footnote{KB excerpt: \texttt{kb/excerpts/jimenez2024-swebench.md\#abstract}.}
\citep[\S{}abstract]{jimenez2024-swebench}:

\begin{quote}\itshape
SWE-bench [is] an evaluation framework consisting of 2,294 software
engineering problems drawn from real GitHub issues and corresponding
pull requests across 12 popular Python repositories.
\end{quote}

Each task instance ships a base commit (the codebase state $s_0$),
the issue text ($\rho$), and the pull-request's test diff partitioned
into two regression-test sets%
\footnote{KB excerpt: \texttt{kb/excerpts/jimenez2024-swebench.md\#sec-eval-protocol}.}.
\textsf{FAIL\_TO\_PASS} is the set of tests that fail on the pre-PR
codebase and must pass on the post-patch state $s_T$.
\textsf{PASS\_TO\_PASS} is the set of tests that pass on the pre-PR
codebase and must continue to pass after the patch. The verifier is
the conjunction
\[
V(s_T, g) = \mathbb{1}\!\left[\bigwedge_{t \in \mathrm{F2P}} \mathrm{pass}(t, s_T)\right] \cdot \mathbb{1}\!\left[\bigwedge_{t \in \mathrm{P2P}} \mathrm{pass}(t, s_T)\right],
\]
which collapses the entire question of ``did the agent fix the bug''
into ``does \texttt{pytest} accept the patch.'' The
\textsf{PASS\_TO\_PASS} guard is what stops trivial regression-by-
deletion: a patch that satisfies the bug-test by breaking unrelated
functionality fails this term \citep[\S{}sec-eval-protocol]{jimenez2024-swebench}.

The launch number is small enough to publish:
``\glsterm{claude}{Claude} 2 \ldots\ resolves only 1.96\% of the issues''
\citep[\S{}sec-headline]{jimenez2024-swebench}. By mid-2025 the
\glsterm{swe-bench-verified}{SWE-Bench Verified} leaderboard sits above $70\%$ for frontier reasoning
agents, a $\sim 35\times$ jump in $\sim 18$ months. Two methodological
artifacts shaped that arc. First, \citet{jimenez2024-swebench} report
that many of the original $2{,}294$ items are under-specified: the
issue text does not uniquely determine which tests count as ``bug
fixed,'' so a correct patch may still fail \textsf{FAIL\_TO\_PASS}.
\glsterm{swe-bench-verified}{SWE-Bench Verified} (OpenAI, August 2024) filters the pool to $500$
human-validated instances where professional software engineers
confirmed unambiguous specs, fair test sets, and reproducible
environments; about one-third of original tasks are rejected at this
filter. Verified is the de-facto leaderboard target as of 2026;
raw SWE-Bench is treated as developer-internal.

\begin{intuition}
SWE-Bench's verifier $V$ is the only one in this section that comes
\emph{free with the source domain}: software engineers already write
regression tests, and the test diff in the canonical PR is the gold
standard the human reviewer used. Every other agentic verifier ---
OS state checks, string-match on a final answer field, DB-row
equality --- is constructed by the benchmark author. This makes
SWE-Bench's score the most directly tied to a real engineering
artifact, and also the most exposed to the spec-ambiguity tail that
Verified was built to cut. (Returning to canonical form: the verifier
is still $V(s_T, g) \in \{0, 1\}$; what differs is that $g$ here is
authored by the original developer, not by the eval team.)
\end{intuition}

The second artifact is the scaffold-vs-model conflation. Most of the
$1.96\% \to 70\%$ jump came from agent harnesses --- SWE-Agent,
OpenHands, Aider, Devin-style scaffolds --- wrapping the same base
model with retrieval, tool loops, and self-reflection cycles, not
from frontier-model capability alone. The headline ``\% resolved'' is
the joint score of model and scaffold, but most leaderboards do not
require scaffold-source disclosure as a leaderboard-entry condition.

\subsection{OSWorld: GUI-grounded computer use}
\label{sec:agentic-osworld}

\citet{xie2024-osworld} push the action space all the way down to the
desktop GUI%
\footnote{KB excerpt: \texttt{kb/excerpts/xie2024-osworld.md\#abstract}.}
\citep[\S{}abstract]{xie2024-osworld}:

\begin{quote}\itshape
We introduce OSWorld, the first-of-its-kind scalable, real computer
environment for multimodal agents, supporting task setup, execution-based
evaluation, and interactive learning across various operating systems
such as Ubuntu, Windows, and macOS.
\end{quote}

The benchmark ships $369$ tasks covering real web and desktop apps,
OS file I/O, and workflows that span multiple applications%
\footnote{KB note: \texttt{kb/notes/evaluation/agentic-benchmarks.md\#23-osworld}.}
\citep[\S{}abstract]{xie2024-osworld}. Each task carries an OS-image
setup script that produces $s_0$, a natural-language task description
$\rho$, and a per-task evaluator script that checks the final OS
state $s_T$ against the goal. The action space is mouse and keyboard
events on the live GUI (typically over VNC), not a structured API:
the agent must convert a natural-language instruction like ``click
the Save button'' into pixel coordinates, which makes vision-based
\glsterm{gui-grounding}{GUI grounding} an unavoidable sub-capability of every score.

The launch headline is the methodological one
\citep[\S{}abstract]{xie2024-osworld}:

\begin{quote}\itshape
While humans can accomplish over 72.36\% of the tasks, the best model
achieves only 12.24\% success, primarily struggling with \glsterm{gui-grounding}{GUI grounding}
and operational knowledge.
\end{quote}

The $72.36\% \to 12.24\%$ gap at submission was the largest among
published 2024 agentic benchmarks. Subsequent computer-use systems
(Anthropic's October 2024 ``computer use'' \glsterm{claude}{Claude} release, and
follow-on agentic VLMs) push the AI number higher, but the failure
mode the abstract names --- \glsterm{gui-grounding}{GUI grounding}, not high-level planning
--- has not gone away.

\begin{contradiction}
OSWorld's metric is bottlenecked by a sub-capability that may not be
the capability the eval was constructed to measure. If a text-only
agent with perfect planning is graded at $0\%$ because it cannot
emit pixel coordinates, the $12.24\%$ baseline is partially measuring
vision-language alignment, not autonomous task completion. The Phase
1 sweep flags this as the ``headless track'' open question: should
benchmarks abstract the GUI to a text representation
(accessibility-tree or DOM) so that the planning-vs-grounding axes
can be measured separately%
\footnote{KB note: \texttt{kb/notes/evaluation/agentic-benchmarks.md\#5-frontier}.}?
The field has not committed.
\end{contradiction}

\subsection{GAIA: general AI assistant}
\label{sec:agentic-gaia}

\citet{mialon2023-gaia} invert the usual benchmark gradient%
\footnote{KB excerpt: \texttt{kb/excerpts/mialon2023-gaia.md\#abstract}.}
\citep[\S{}abstract]{mialon2023-gaia}:

\begin{quote}\itshape
We propose GAIA, a benchmark for General AI Assistants \ldots GAIA
proposes real-world questions that require a set of fundamental
abilities such as reasoning, multi-modality handling, web browsing,
and generally tool-use proficiency. GAIA questions are conceptually
simple for humans yet challenging for most advanced AIs.
\end{quote}

The headline gap is the load-bearing one: ``human respondents obtain
$92\%$ vs.\ $15\%$ for GPT-4 equipped with plugins''
\citep[\S{}abstract]{mialon2023-gaia}. The $92\% \to 15\%$ inversion
is the design point: every GAIA question is solvable by an average
human, and the difficulty for the model is in the multi-step nature
of the tool chain, not in specialized knowledge.

The dataset is $466$ questions total, with $300$ retained for the
gated leaderboard test set and $166$ released as a public development
set, stratified into Levels 1/2/3 by tool-chain depth%
\footnote{KB excerpt: \texttt{kb/excerpts/mialon2023-gaia.md\#sec-dataset}.}.
Level 1 is one-to-three-tool short chains with no specialized
knowledge; Level 2 is multi-tool web-browsing plans; Level 3 is
long-horizon tool use with multi-document synthesis, sometimes more
than ten hops. Answers are checked by exact-match-after-normalization
on a single short string --- a number, a date, a name --- not by
\glsterm{llm-as-judge}{LLM-as-judge}. The verifier is mechanical and reproducible, which
matters because the test gold answers are kept private and submissions
are run server-side by the GAIA team to delay contamination via
training-data leakage. The $92\%$ expert baseline and the $15\%$
frontier-launch number set up exactly the diagnostic the abstract
calls out: the gap is in robustness over multi-step plans, not in
any single capability the model lacks.

The 2026 frontier picture: leaderboard entries above $50\%$ now exist,
narrowing but not closing the human gap on Levels 1 and 2; Level 3
remains far from \glsterm{saturation}{saturation}. Why expert humans hit $\sim 92\%$ and
frontier models hit $\sim 52\%$ at the time of writing is mostly
the long-horizon failure mode --- accumulating small mistakes over
ten-plus tool calls compounds into wrong final-string answers --- not
single-step capability gaps.

\subsection{$\tau$-Bench: tool-using agent in conversation}
\label{sec:agentic-tau-bench}

\citet{yao2024-tau-bench} add the two pieces missing from the others:
a simulated user who can clarify and push back, and a policy document
the agent must follow%
\footnote{KB excerpt: \texttt{kb/excerpts/yao2024-tau-bench.md\#abstract}.}
\citep[\S{}abstract]{yao2024-tau-bench}:

\begin{quote}\itshape
Existing benchmarks do not test language agents on their interaction
with human users or ability to follow domain-specific rules, both of
which are vital for deploying them in real world applications.
\end{quote}

The benchmark construction%
\footnote{KB excerpt: \texttt{kb/excerpts/yao2024-tau-bench.md\#sec-user-sim}.}
\citep[\S{}sec-user-sim]{yao2024-tau-bench}:

\begin{quote}\itshape
The benchmark employs simulated conversations between a user
(generated by language models) and an agent equipped with
domain-specific API tools and policy guidelines.
\end{quote}

Two domains ship in the original release: \emph{retail} (e-commerce
return-and-exchange policy) and \emph{airline} (rebooking and
policy enforcement). The verifier is end-state DB equality --- the
final database rows must match the annotated goal state ---
sidestepping \glsterm{llm-as-judge}{LLM-as-judge} contamination on the agent side. The
methodologically novel piece is the $\mathrm{pass}^k$ metric%
\footnote{KB excerpt: \texttt{kb/excerpts/yao2024-tau-bench.md\#sec-passk}.}
\citep[\S{}sec-passk]{yao2024-tau-bench}:
\[
\mathrm{pass}^k(\text{agent}, \text{task}) = \mathbb{1}\!\left[\bigwedge_{i=1}^{k} V(s_T^{(i)}, g) = 1\right],
\]
the indicator that \emph{all} $k$ independent runs succeed --- the
dual of the standard $\mathrm{pass}@k$ (``at least one of $k$
succeeds''). The two metrics measure orthogonal properties.
$\mathrm{pass}@k$ measures whether the capability exists somewhere in
the sample distribution; $\mathrm{pass}^k$ measures whether the
capability is reliable enough to deploy without human review.
The launch finding is the gap between them
\citep[\S{}sec-result]{yao2024-tau-bench}:

\begin{quote}\itshape
even state-of-the-art function calling agents (like gpt-4o) succeed
on $<50\%$ of the tasks, and are quite inconsistent ($\mathrm{pass}^8 < 25\%$
in retail).
\end{quote}

A $50\%$ $\mathrm{pass}@1$ agent with $25\%$ $\mathrm{pass}^8$ has
roughly half its successes non-reproducible under user-simulator
variation. For production customer-service deployment this is the
load-bearing quantity.

\begin{intuition}
$\mathrm{pass}^k$ is the production discipline of agent evaluation.
The same $50\%$ headline number on $\mathrm{pass}@1$ can mean two
operationally opposite things: (i) the agent is reliably right on
half the tasks and reliably wrong on the other half (high
$\mathrm{pass}^k$ on the right half), or (ii) the agent is
unreliably right on every task with a flip-of-the-coin success
distribution (low $\mathrm{pass}^k$ everywhere). For a customer-service
system, the first failure mode is repairable by routing rules and
the second is not. (Returning to canonical form: $V(s_T^{(i)}, g)$ is
the same indicator in both cases; the difference is variance over
trajectories.)
\end{intuition}

\subsection{Realism vs measurability}
\label{sec:agentic-realism-vs-measurability}

The four benchmarks make a single tradeoff explicit:
\emph{the more an environment looks like real deployment, the harder
it is to score reliably}. \glsterm{swe-bench-verified}{SWE-Bench Verified} is the most ecologically
realistic --- the codebases are real, the issues are real, the tests
are the developers' own --- but its construction history (the original
$2{,}294$-task set being filtered down to $500$ human-vetted
instances) records exactly how much spec ambiguity, environment
non-reproducibility, and unfair test sets the realism cost.
OSWorld is realistic at the OS-image level but the metric absorbs
vision-grounding noise that confounds the autonomous-task-completion
construct it advertises. GAIA buys very-low contamination and a
mechanical verifier by hand-crafting items, but the cost is item
scarcity ($466$ total) and the slow gated-evaluation cycle. $\tau$-Bench
is the most measurable --- DB-row equality is unambiguous --- but
the user side of the conversation is itself an \glsterm{llm}{LLM}, which makes the
score partly a measurement of how well the agent handles the
specific generation distribution of the user simulator, not a
deployment-grade user interaction.

The same tradeoff is present, on a different axis, in
\cref{sec:knowledge-benchmarks}: AIME-style benchmarks have
unambiguous integer-answer grading and limited realism; SWE-Bench
Verified has high realism and subjective items in the long tail
that survived the human-validation filter. The agentic frontier and
the knowledge frontier face the same problem from opposite directions:
adding realism to the eval surface costs grading reliability, and
the field has not yet found a construction that does not pay this
cost.

\subsection{What slips through}
\label{sec:agentic-slips-through}

The \glsterm{end-state-verifier}{end-state verifier} $V(s_T, g)$ is silent on a class of failure
modes that show up only in the trajectory $(s_0, a_0, s_1, \ldots, s_T)$.
Three failure modes are now documented enough to name. \emph{Long-horizon
coherence loss} --- the agent succeeds on Level 1 GAIA tasks and
collapses on Level 3, even though every intermediate Level-2 step is
in its capability surface --- is observable in GAIA's level-stratified
scores but not in the aggregate accuracy a model card would report.
\emph{Deception of the user simulator} in $\tau$-Bench --- the agent
satisfies the simulator and the policy document but reaches the goal
DB state by misrepresenting the situation to the simulated user ---
is exactly the kind of behavior the policy-compliance check was
designed to expose, but the DB-equality verifier rewards it. And
\emph{skill memorization} on SWE-Bench --- the issue thread, the
canonical fix, and even the test diff are public on GitHub; nothing
in the verifier $V$ distinguishes ``the model reproduced the
canonical patch from training-data memory'' from ``the model
reasoned about the bug from first principles.''

\begin{contradiction}
Are agentic benchmarks now the bottleneck on capability evaluation,
or are they leaky proxies that overstate generality? Both readings
are defensible from the 2026 evidence. The bottleneck reading: the
$1.96\% \to 70\%$ \glsterm{swe-bench-verified}{SWE-Bench Verified} arc and the $15\% \to 50\%+$
GAIA arc are the most informative single capability measurements
the field has produced since \glsterm{mmlu}{MMLU} saturated, and the realism gap is
closing fast enough that the next eval generation will need to
target what current agentic benchmarks miss. The leaky-proxy reading:
SWE-Bench scores are conflated with scaffold engineering, OSWorld
scores are bottlenecked on vision grounding, GAIA's privacy of gold
answers cannot be audited externally, and $\tau$-Bench's user
simulator is itself a model in training distribution. The field's
own response --- replacement benchmarks (Multi-SWE-bench, SWE-Lancer,
$\tau$2-Bench), cost-aware leaderboards (Aider Polyglot's
$\$$-per-task axis), and \glsterm{task-horizon-metric}{METR}'s \glsterm{task-horizon-metric}{task-horizon metric}%
\footnote{KB note: \texttt{kb/notes/evaluation/agentic-benchmarks.md\#5-frontier}.}
--- is consistent with believing both readings simultaneously: agentic
benchmarks are the best capability instruments we have, and we do
not yet trust them to predict deployment \glsterm{value}{value}.
\end{contradiction}

The next layer of the layered defense, taken up in
\cref{sec:safety-adversarial-robustness}, asks not what the agent
\emph{can} do but what it \emph{will} do under adversarial pressure
--- the jailbreak-resistance and \glsterm{dangerous-capability-evaluation-2}{dangerous-capability evaluation}
surface, where the construct is no longer task completion but
behavior under deliberate attack.
